% Generated draft: requires mathematical and source review.
\documentclass[11pt]{article}
\usepackage[margin=25mm]{geometry}
\usepackage{fontspec}
\setmainfont{Latin Modern Roman}
\usepackage{amsmath,amssymb,mathtools,unicode-math}
\setmathfont{Latin Modern Math}
\usepackage{enumitem,needspace,xurl,hyperref}
\setlength{\parindent}{0pt}
\setlength{\parskip}{5pt}
\setlength{\emergencystretch}{3em}
\newcommand{\N}{\mathbb{N}}
\newcommand{\Z}{\mathbb{Z}}
\newcommand{\Q}{\mathbb{Q}}
\newcommand{\R}{\mathbb{R}}
\newcommand{\C}{\mathbb{C}}
\newcommand{\F}{\mathbb{F}}
\newcommand{\A}{\mathbb{A}}
\newcommand{\PP}{\mathbb{P}}
\newcommand{\E}{\mathbb{E}}
\newcommand{\eps}{\varepsilon}
\newcommand{\dd}{\,\mathrm{d}}
\begin{document}
\section*{UnsolvedMath ID: 30006442}
\textbf{Source problem code:} catalogue entry
\subsection{Definitions and mathematical statement}
The supplied cleaned formulation concerns a set \(X\subseteq\mathbb{C}_p^n\) definable in a Hensel-minimal structure, where \(\mathbb{C}_p\) conventionally denotes the completion of an algebraic closure of \(\mathbb{Q}_p\). The record does not specify the prime \(p\), the dimension restrictions on \(n\), the structure, the permitted parameters, or the precise Hensel-minimality hypothesis.

The notation \(X^{\mathrm{trans}}(\mathbb{Q},H)\) is not defined in the supplied excerpt. To display the requested quantifiers without inventing those definitions, let \(T_X\subseteq X\) denote the intended transcendental part and let \(h\) denote the intended height on rational points, both still to be specified from the source. Write, schematically,
\[
N_{X,T_X,h}(H)=\#\{x\in T_X\cap\mathbb{Q}^n:h(x)\leq H\},\qquad H\in\mathbb{R},\ H\geq1,
\]
using the usual inclusion \(\mathbb{Q}\hookrightarrow\mathbb{C}_p\) in each coordinate. This notation is a placeholder for the missing source definitions, not a choice of transcendental part or height.

\textbf{Question as supplied, with unresolved qualifications.} Can the number \(\#X^{\mathrm{trans}}(\mathbb{Q},H)\) be bounded uniformly for all \(H\geq1\), with upper bounds as in the source's Theorem~2? On the literal interpretation of uniform boundedness in \(H\), the schematic assertion for each admissible \(X\) is
\[
\exists B<\infty\quad\forall H\in\mathbb{R}_{\geq1},\qquad N_{X,T_X,h}(H)\leq B.
\]
Here independence from \(H\) is the only uniformity explicitly requested; the allowed dependence of \(B\) on \(X\), \(p\), the structure, parameters, or complexity data is unspecified. The source's Theorem~2 and reference~[15], said to treat the one-dimensional case, are absent from the supplied record. Consequently, this schematic assertion cannot be identified confidently with the full intended theorem-level question: neither the required form of the bounds nor their dependence on auxiliary data is available. The supplied cleaned statement paraphrases the requested comparison as bounds analogous to those known in dimension one. A complete precise formulation requires the omitted definitions and theorem.
\subsection{Short English statement}
The supplied question asks whether the transcendental rational points of a Hensel-minimally definable set admit a bound uniform in the height cutoff \(H\geq1\), with the quantitative features of an omitted Theorem~2. The transcendental part, height convention, hypotheses, and allowed dependence of the bound remain unspecified. The record labels the problem open; this is supplied dated curation, not an independently verified status.
\subsection{Sources}
\begin{itemize}
\item Local UnsolvedMath registry, permanent ID \(30006442\), title: \emph{Uniform Bounds for Transcendental Rational Points}; source problem code \texttt{catalogue entry}. The supplied record belongs to the Oberwolfach Reports Open Problems collection. No individual proposer is identified; the recorded proposed year is 2025.
\item Supplied report citation: \emph{Singularities (2026)}, Oberwolfach Reports, DOI \texttt{10.4171/owr/2025/45}, pp.~2413--2490. \url{https://doi.org/10.4171/owr/2025/45}. The literature-source list labels this same reference \emph{Oberwolfach Report 45/2025}. The discrepancy between the displayed year 2026 and the 2025 report identifiers is retained without resolution.
\item Supplied literature reference: \emph{Effective Pila--Wilkie bounds}. \url{https://www.wisdom.weizmann.ac.il/~dnovikov/Papers/Pfaff_Wilkie.pdf}. No authors or full bibliographic details for this reference are supplied.
\item The original excerpt attributes its inspiration to reference~[15], saying that the one-dimensional case is treated there, and requests upper bounds as in Theorem~2. Neither reference~[15] nor Theorem~2 is supplied, and their relationship to the separately listed literature reference is not established.
\item Literature assessment dated exactly \texttt{2026-08-22}: the supplied curation labels the question open and states that existing Pila--Wilkie-type estimates do not settle the requested uniform bound with the desired theorem-level constants. This is supplied dated curation, not independently verified. The background explicitly calls for independent verification of time-sensitive status and cited claims.
\item The registry labels the statement \texttt{corrected\_verified} and reports that the official report was checked, the question isolated, and notation repaired. These are supplied verification claims; this draft does not independently verify them or claim to have visited either linked source.
\end{itemize}
\end{document}
